\section*{Chapter \ref{ch:fm:treasury-and-funding} --- Treasury and Funding}

\begin{solution}{exo:fm:treasury-and-funding:1}
A: $0.04\times120+0.25\times120+0.15\times60=43.8$; B: $0.10\times120=12.0$; C: $0.05\times120+0.15\times120+0.30\times80=48.0$ (\$ million).
\end{solution}

\begin{solution}{exo:fm:treasury-and-funding:2}
\$5 million (50\%) on the long; 150\% of the short's value, \$15 million, including the \$10 million proceeds of the sale.
\end{solution}

\begin{solution}{exo:fm:treasury-and-funding:3}
$38.1\times5\%=\$1.91$ million a year.
\end{solution}

\begin{solution}{exo:fm:treasury-and-funding:4}
Setting $x_{ib}=1$ for one broker and 0 elsewhere satisfies $\sum_bx_{ib}=1$ and the bounds, with the auxiliary variables at their smallest
feasible values; the programme minimises over a set containing that point, so its minimum cannot exceed that point's cost.
\end{solution}

\begin{solution}{exo:fm:treasury-and-funding:5}
Its margin, \$72.0 million, is almost A's \$71.55 million, but its financing spreads are the lowest: \$2.39 million a year against A's \$3.30
million.
\end{solution}

\begin{solution}{exo:fm:treasury-and-funding:6}
Ten days postpone the day the cash runs out from day 5 to day 11, just beyond the brokers' increases in the scenario's first week, but the increases
then arrive; thirty days cover the whole stress and the firm never runs out.
\end{solution}

\begin{solution}{exo:fm:treasury-and-funding:7}
With \$50 or \$60 million the optimised book lasts the month; the book at broker B runs out on day 4 with \$50 million and on day 5 with \$60
million.
\end{solution}

\begin{solution}{exo:fm:treasury-and-funding:8}
Lower margin today is not more liquidity in stress: a single broker can raise its \omterm{def:fm:treasury-and-funding:house}{house margin} at will, and with no second broker and no lock-up
the whole book's requirement moves at once. In the chapter the book at one broker runs out two days earlier than the diversified book and needs
\$29.8 million more of buffer.
\end{solution}

\begin{solution}{pb:fm:treasury-and-funding:1}
\begin{enumerate}
\item See \cref{def:fm:treasury-and-funding:cash}.
\item Initial margin at clearing houses rose by roughly \$300 billion over the month, and variation margin calls were far above February's.
\item See \cref{def:fm:treasury-and-funding:house,def:fm:treasury-and-funding:lockup}.
\item Regulation T: 50\% initial margin on margin equity securities, 150\% on short sales; FINRA 4210: 25\% maintenance on longs, or portfolio
margin.
\item See \cref{def:fm:treasury-and-funding:optim}.
\item See \cref{prop:fm:treasury-and-funding:lp}.
\item See \cref{tab:fm:treasury-and-funding:alloc}: from \$33.9 million and \$4.98 million a year optimised to \$144.0 million and \$10.92 million
all at C.
\item It balances each broker's slice to avoid A's net charge and C's add-ons, and sends B the positions its flat gross rate prices cheaply.
\item Relationships, repricing, the cost of moving, and the value of a broker that supports the firm in stress.
\item See \cref{def:fm:treasury-and-funding:horizon}.
\item Five days of longs $-2\%$, shorts $-1\%$, index $-2\%$; brokers double \omterm{def:fm:treasury-and-funding:house}{house margin} by day 5. The week loses \$15.5 million.
\item \$43.0 and \$71.8 million without lock-up; about \$36 and \$37 million with; the clearing broker's margin \$2.02 and \$2.13 million.
\item Optimised: day 5, buffer \$8.1 million more; one broker: day 3, \$29.8 million more; with a thirty-day lock-up, no shortfall.
\item It postpones the shortfall to the end of the notice period (day 6 for five days, day 11 for ten).
\item See \cref{dat:fm:treasury-and-funding:fsb}.
\item By \cref{met:fm:treasury-and-funding:buffer}: the largest cumulative shortfall over the stress scenarios, over the longest notice period.
\item The lock-up: the optimisation saves \$1 million a year and shortens nothing; without a lock-up the optimised book still runs out on day 5.
\item Lock-ups of at least the length of the stress for the house terms that move most; release of excess on demand; notice before any change
of eligible collateral.
\item \$38.1 million (53\%) of margin saved against the cheapest single broker; without a lock-up the cash lasts to day 5 (day 3 at one broker),
with a thirty-day lock-up the whole month.
\item Hold the book across several brokers at the least cost that remains acceptable when terms change, with lock-ups covering the stress horizon;
keep a buffer equal to the largest shortfall of the stress forecast, reported monthly with the \omterm{def:fm:treasury-and-funding:horizon}{survival horizon}.
\end{enumerate}
\end{solution}

\begin{solution}{iq:fm:treasury-and-funding:1}
The margin a broker requires under its own rules; it prices the broker's risk on the client (concentration, liquidity, the client's credit) beyond
what the regulatory floor assumes.
\iqlookfor{the broker's own risk and discretion.}
\end{solution}

\begin{solution}{iq:fm:treasury-and-funding:2}
The days until \omterm{def:fm:treasury-and-funding:cash}{unencumbered cash}, after all calls and outflows in a stress, first goes negative; from a daily forecast of every account's P\&L and
requirement.
\iqlookfor{a daily forecast, not a ratio.}
\end{solution}

\begin{solution}{iq:fm:treasury-and-funding:3}
Minimise the cost of margin plus financing over fractions $x_{ib}$ with $\sum_bx_{ib}=1$; linearise the absolute net and the concentration excesses
with auxiliary variables; a linear programme.
\iqlookfor{convexity and the linearisation.}
\end{solution}

\begin{solution}{iq:fm:treasury-and-funding:4}
Whether the agreement allows it without notice, the new requirement against \omterm{def:fm:treasury-and-funding:cash}{unencumbered cash}, what can move to other brokers, and which positions
to reduce; then pay, reduce or move, and report the \omterm{def:fm:treasury-and-funding:horizon}{survival horizon}.
\iqlookfor{terms, cash and options, in that order.}
\end{solution}

\begin{solution}{iq:fm:treasury-and-funding:5}
To keep several brokers engaged, to buy a lock-up or committed financing, and to keep positions where they can be moved or financed in stress.
\iqlookfor{the price of resilience.}
\end{solution}

\begin{solution}{iq:fm:treasury-and-funding:6}
Map the accounts and their terms; build scenarios of P\&L, house-margin increases, withheld excess and clearing-margin moves; forecast cash daily;
report the \omterm{def:fm:treasury-and-funding:horizon}{survival horizon} and the buffer needed, and the terms that drive them.
\iqlookfor{joint scenarios and a daily forecast.}
\end{solution}
